\section{Model and exact growth theorem}
\label{van:section}

The local redirection models of Hartle, Krapivsky, Redner, and Zhang
redirect a core with positive leaf count to one of its leaves
\cite{HRKZ2026}. The resulting independence of different core families
is the basis of the proof below.

Fix a finite parameter $\omega\ge0$.  A \emph{core} is a vertex of degree
at least two, and its \emph{leaf count} is the number of its neighbours
of degree one.  In continuous time the rules are:
\begin{enumerate}
\item Each leaf has a rate-one clock.  When it rings, attach a new vertex
to that leaf's unique neighbour.
\item Each core with positive leaf count has a rate-one clock.  When it
rings, choose one of its leaves uniformly and attach a new vertex to
that leaf.  The chosen leaf becomes a new core.
\item Each core with zero leaves has a rate-$\omega$ clock.  When it
rings, attach a new vertex directly to that core.
\end{enumerate}
The embedded tree at successive vertex births is precisely
$\operatorname{VAN}(\omega)$: $\omega=0$ gives PAN, after deleting
rejected selections, and $\omega=1$ gives DAN.  These rules are understood
after the initial seed has at least three vertices.  A two-vertex seed
reaches a three-vertex path after its first event and is included by
this finite initial delay.

Write $C(t)$, $L(t)$ and $C_0(t)$ for the numbers of cores, leaves and
zero-leaf cores.  Then $N(t)=C(t)+L(t)$ and the total vertex-birth rate is
\[
 L(t)+C(t)-C_0(t)+\omega C_0(t)
 =N(t)+(\omega-1)C_0(t)\le \max\{1,\omega\}N(t).
\]
In particular, the process is nonexplosive.  Let $C_n$ be the core count
when the tree has $n$ vertices.

\begin{theorem}[Exact modified-model exponent and almost-sure amplitude]
\label{van:main}
For every finite $\omega\ge0$, define, for $s>0$,
\begin{align}
 F(s)&=\int_0^1 x^s e^{x-1}\,dx,
 \label{van:F}\\
 m_\omega(s)&=
 \frac{(s+\omega)(1-F(s))}
 {s\bigl(s+\omega-\omega F(s)\bigr)}.
 \label{van:m}
\end{align}
There is a unique $\alpha_\omega\in(0,1)$ satisfying
$m_\omega(\alpha_\omega)=1$.  Equivalently,
\begin{equation}
 \boxed{\quad
 F(\alpha_\omega)=
 \frac{(1-\alpha_\omega)(\alpha_\omega+\omega)}
 {\omega+(1-\omega)\alpha_\omega}.
 \quad}
 \label{van:root}
\end{equation}
For every finite tree seed with at least two vertices, there are random
variables $A,V\in(0,\infty)$ such that, almost surely,
\begin{equation}
 e^{-\alpha_\omega t}C(t)\longrightarrow A,
 \qquad e^{-t}N(t)\longrightarrow V,
 \qquad
 \frac{C_n}{n^{\alpha_\omega}}\longrightarrow A V^{-\alpha_\omega}.
 \label{van:limits}
\end{equation}
The exponent $\alpha_\omega$ is effectively computable whenever
$\omega$ is computable.  More quantitatively,
\begin{equation}
 -m_\omega'(s)\ge\frac1{(1+s)^2}\ge\frac14,
 \qquad 0<s<1.
 \label{van:slope}
\end{equation}
\end{theorem}

The characteristic equation~\eqref{van:root} is the equation derived
in \cite[Eq.~(25)]{HRKZ2026}. The theorem identifies its unique root
with the almost-sure exponent and proves the positive-amplitude
limit for every finite seed. The parameters $\alpha_0$ and $\alpha_1$
consequently satisfy
$F(\alpha_0)=1-\alpha_0$ and $F(\alpha_1)=1-\alpha_1^2$.
Throughout the theorem, $\omega$ is finite.  Its assertion does not
include the separate limiting model $\operatorname{VAN}(\infty)$.

\subsection{The exact life process of a core}

Regard a newly created core as an individual whose children are the
cores it subsequently creates.  Every such newborn individual starts
with one leaf.  Its own leaf count $X$ has transition rates
\begin{equation}
 \ell\longrightarrow\ell+1 \text{ at rate }\ell,\qquad
 \ell\longrightarrow\ell-1 \text{ at rate }1
 \quad(\ell\ge1),\qquad
 0\longrightarrow1 \text{ at rate }\omega.
 \label{van:life}
\end{equation}
Each downward transition creates one child, itself starting at leaf
count one.  The parent loses the promoted leaf and the child acquires
the newly attached leaf.

The life processes of different individuals can be constructed
independently.  Their transition rates do not involve their core
neighbours; fresh future clock increments at a promotion are
independent of the past.  Thus this is a genuine single-type
Crump--Mode--Jagers process.  The exterior core edge of a newborn
individual does not affect its life process.  This observation also
allows us to study one genealogical family in isolation.

\begin{lemma}[An exact excursion transform]
\label{van:excursion}
For the chain \eqref{van:life} started at $\ell\ge1$, let $T_0$ be its
first hit of zero, prior to any restart from zero.  For $s>0$,
\begin{equation}
 h_\ell(s):=\mathbb E_\ell[e^{-sT_0};T_0<\infty]
 =\frac1{(\ell-1)!}\int_0^1
 x^s(1-x)^{\ell-1}e^{x-1}\,dx.
 \label{van:hitting}
\end{equation}
In particular, $h_1(s)=F(s)$.
\end{lemma}

\begin{proof}
Set $h_0=1$ and define $h_\ell$ by the integral on the right.
Integration of the derivative of
$x^{s+1}(1-x)^{\ell-1}e^{x-1}$ gives
\[
 (s+\ell+1)h_\ell=\ell h_{\ell+1}+h_{\ell-1},
 \qquad \ell\ge1.
\]
For $\ell=1$ the boundary value at $x=1$ is one and supplies the
term $h_0$; for $\ell\ge2$ both boundary values vanish.
Moreover $0\le h_\ell\le1/\ell!$.
The displayed recurrence is the discounted backward equation.
Apply its bounded stopped-martingale identity up to $T_0$, time $t$,
and the first hit of level $R$.  The contribution at level $R$
vanishes as $R\to\infty$, since $h_R\to0$; nonexplosion justifies
removing this localization.  The remaining time-$t$ term is bounded
by $e^{-st}$ and vanishes as $t\to\infty$.  This proves
\eqref{van:hitting}.
\end{proof}

Let $\xi$ be the reproduction point process of an individual started
at one, and let $\nu_\omega=\mathbb E\xi$ be its mean reproduction
measure.  The downward-event intensity gives
\begin{equation}
 \widehat\nu_\omega(s)
 =\mathbb E\int_0^\infty e^{-st}\xi(dt)
 =\mathbb E_1\int_0^\infty e^{-st}\mathbf1_{\{X_t>0\}}\,dt.
 \label{van:reproduction}
\end{equation}
Before $T_0$ this expected occupation is $(1-F(s))/s$.
Upon hitting zero, the chain waits an independent
$\operatorname{Exp}(\omega)$ time before restarting from one.
For $\omega=0$ this waiting time is infinite, with discounted
transform zero.  The strong Markov property therefore gives
\[
 \widehat\nu_\omega(s)
 =\frac{1-F(s)}s+
 F(s)\frac{\omega}{s+\omega}\widehat\nu_\omega(s).
\]
Solving proves $\widehat\nu_\omega(s)=m_\omega(s)$ as in
\eqref{van:m}; all its displayed denominators are positive for $s>0$.

The first child is born at age $B\sim\operatorname{Exp}(1)$.
Indeed, before the first core-clock ring only upward leaf moves
can occur, so the core remains active.  More generally, the entire
reproduction process is obtained by accepting a subset of the points
of a rate-one Poisson process $\Pi$, rejecting its points when $X=0$.
Consequently
\begin{align}
 m_\omega(s)&\le \frac1s,
 &
 -m_\omega'(s)&\ge\mathbb E[B e^{-sB}]
                  =\frac1{(1+s)^2},
 \label{van:first-child}\\
 \mathbb E\left[\left(\int_0^\infty e^{-\theta t}\xi(dt)\right)^2\right]
 &\le \frac1{\theta^2}+\frac1{2\theta},
 &&\theta>0.
 \label{van:second-moment}
\end{align}
There is positive probability of reaching zero and spending positive
time there, for every finite $\omega$.  Hence $m_\omega(1)<1$.
Since $F(0)=1-e^{-1}$, formula \eqref{van:m} also shows
$m_\omega(s)\to\infty$ as $s\downarrow0$.  Strict decrease and continuity
now give the unique root in $(0,1)$.  Algebraic rearrangement gives
\eqref{van:root}; its denominator
$\omega+(1-\omega)s=s+\omega(1-s)$ is positive on that interval.

\subsection{Branching growth and arbitrary finite seeds}

\begin{lemma}[Ordinary-family growth and a continuous-time mean bound]
\label{van:ordinary}
Put $\alpha=\alpha_\omega$ and
$\beta=-m_\omega'(\alpha)\in(0,\infty)$.
For one ordinary family started with a core of leaf count one,
\begin{equation}
 e^{-\alpha t}C^{(1)}(t)\longrightarrow
 A_1=\frac{W}{\alpha\beta}\in(0,\infty)
 \quad\text{almost surely}.
 \label{van:ordinary-limit}
\end{equation}
Here $W$ is the limiting intrinsic CMJ martingale.
For every $\sigma>\alpha$,
\begin{equation}
 \mathbb E C^{(1)}(t)
 \le \frac{e^{\sigma t}}{1-m_\omega(\sigma)}.
 \label{van:mean-bound}
\end{equation}
\end{lemma}

\begin{proof}
We verify the conditions for the CMJ law of large numbers
\cite[Section~2.2, equation~(2.13) and its almost-sure
upgrade]{IKM2024}.
Reproduction processes are independent copies by the exact
construction above.  Their mean measure is nonlattice, since the
first birth age is exponentially distributed.  The Malthusian
equation is $m_\omega(\alpha)=1$, and
\[
 0<\frac1{(1+\alpha)^2}\le\beta\le
 \int_0^\infty t e^{-\alpha t}\,dt=\alpha^{-2}<\infty.
\]
The moment bound \eqref{van:second-moment} holds for every
$\theta>0$, including a $\theta<\alpha/2$.
It supplies both the discounted $x\log x$ condition
(see \cite{Olofsson1998}) and the stronger
moment hypothesis in \cite[Remark~2.1]{IKM2024}.  The counting
characteristic $\phi(t)=\mathbf1_{\{t\ge0\}}$ is bounded and has the
required exponentially weighted regularity and summability.
For this characteristic, the weighted local suprema are bounded by
a constant times $e^{-\alpha t}\mathbf1_{\{t\ge-1\}}$, so the
characteristic condition in the almost-sure upgrade is satisfied.
Nerman's theorem gives the limit
$\beta^{-1}\int_0^\infty e^{-\alpha t}\,dt\;W
=W/(\alpha\beta)$ almost surely.

We give an explicit nondegeneracy check that also permits infinitely
many lifetime offspring. Let $\sigma_v$ be the birth time of an
individual $v$, and set
\[
 W_j=\sum_{|v|=j}e^{-\alpha\sigma_v},\qquad
 Y=\int_0^\infty e^{-\alpha t}\xi(dt).
\]
The generation martingale $(W_j)$ has mean one and the same limit $W$
as the intrinsic CMJ martingale
\cite[Section~2.2, equation~(2.12)]{IKM2024}. Independence of the child
families and monotone truncation of the discounted offspring sum give
\[
 \operatorname{Var}(W_{j+1})
 =m_\omega(2\alpha)\operatorname{Var}(W_j)+\operatorname{Var}(Y).
\]
Here $m_\omega(2\alpha)<1$ and $Y$ has finite second moment by
\eqref{van:second-moment}. Thus $W_j$ is bounded in $L^2$, so
$\mathbb E W=1$. The branching identity for $W$ implies
$q=\mathbb E[q^K]$, where $q=\mathbb P(W=0)$ and $K$ is the
possibly infinite number of children of one individual. We have
$K\ge1$ almost surely. Moreover $\mathbb P(K\ge2)\ge1/2$, since
an upward first move (probability $1/2$) leaves the core active
after its first child is born. As $q<1$, the identity forces $q=0$.
This proves $W>0$ almost surely. Local finiteness holds because the
total core-birth rate is at most the core count, giving a Yule upper
coupling.

For the mean bound, set $M(t)=\mathbb E C^{(1)}(t)$.
The ordinary branching decomposition gives
\[
 M(t)=1+\int_0^t M(t-u)\nu_\omega(du).
\]
After multiplication by $e^{-\sigma t}$ the renewal kernel has
total mass $m_\omega(\sigma)<1$, while the forcing is at most one.
Its nonnegative renewal series is bounded by
$\sum_{j\ge0}m_\omega(\sigma)^j$, which proves
\eqref{van:mean-bound}.
\end{proof}

\begin{lemma}[Positive-event inheritance of the amplitude]
\label{van:inheritance}
Let $C^{(\ell)}(t)$ be the size of the family of one initial core
whose leaf count is $\ell$.  For every integer $\ell\ge1$,
$e^{-\alpha t}C^{(\ell)}(t)$ has a finite strictly positive limit
almost surely.  The same holds for $\ell=0$ when $\omega>0$.
For $\ell=0$ and $\omega=0$ the limit is zero.
\end{lemma}

\begin{proof}
The case $\ell=1$ is Lemma~\ref{van:ordinary}.
For $\ell\ge2$, start with one ordinary core and stop at
\[
 \tau=\min\{\text{first hit of leaf count }\ell,\,
                  \text{first core birth}\}.
\]
The event $E$ that the hit occurs first has probability
\[
 \mathbb P(E)=\prod_{j=1}^{\ell-1}\frac{j}{j+1}=\frac1\ell>0.
\]
On $E$, there is still exactly one core, with leaf count $\ell$,
and no descendants.  The strong Markov future, conditionally on
the history up to $\tau$, is exactly the family started at $\ell$.
The property that $e^{-\alpha t}C(t)$ has a finite strictly positive
limit is preserved by every finite time shift.  It holds almost surely
for the ordinary family, hence also conditionally on $E$.
The Markov property therefore proves that it holds with probability
one for a family started at $\ell$.

If $\ell=0$ and $\omega>0$, an independent finite
$\operatorname{Exp}(\omega)$ activation delay is followed by the
ordinary family.  Its limiting amplitude is the ordinary amplitude
multiplied by the strictly positive discount of this delay.
For $\ell=0$ and $\omega=0$, the family consists of its one inert
ancestor forever.
\end{proof}

For a finite tree with at least three vertices, assign each initial
core its own genealogical family.  Pre-existing core edges do not enter
any rate in \eqref{van:life}, so these are finitely many independent
families with their respective initial leaf counts.  All new cores
belong to exactly one such family.  At least one initial core has
positive leaf count.  Lemma~\ref{van:inheritance} consequently gives
\begin{equation}
 e^{-\alpha t}C(t)\longrightarrow A\in(0,\infty)
 \quad\text{almost surely}.
 \label{van:seed-amplitude}
\end{equation}

We will also need a mean bound for this seed.  The reproduction
intensity of each exceptional initial core is at most one.
Every child begins an ordinary family, so \eqref{van:mean-bound}
and a convolution with this initial reproduction measure show that,
for every $\sigma>\alpha$,
\begin{equation}
 \mathbb E C(t)\le K_\sigma e^{\sigma t}
 \label{van:seed-mean}
\end{equation}
for some finite constant depending on the seed.  For example, if
there are $k$ initial cores, one may take
$K_\sigma=k\bigl(1+[\sigma(1-m_\omega(\sigma))]^{-1}\bigr)$.
This is an auxiliary continuous-time bound; no assertion about
discrete-time mean asymptotics is being used.

\subsection{The exact population clock}

A core birth removes one leaf from its parent and supplies one leaf
to the newborn core.  It leaves the \emph{total} leaf count unchanged.
All other vertex births increase that count by one.  Hence
\begin{equation}
 L(t)\ \text{jumps by }+1\text{ at rate }\ L(t)+\omega C_0(t).
 \label{van:leaf-rate}
\end{equation}
Choose $\alpha<\sigma<1$ in \eqref{van:seed-mean}.
The compensated counting process yields
\begin{equation}
 e^{-t}L(t)=L(0)+\omega\int_0^t e^{-u}C_0(u)\,du+\mathcal M_t,
 \label{van:leaf-martingale}
\end{equation}
where $\mathcal M$ is a square-integrable martingale on every finite
interval.  The drift integral has finite expectation as $t\to\infty$.
Moreover the mean version of \eqref{van:leaf-martingale} gives
\[
 \mathbb E L(t)
 =e^t\left(L(0)+\omega\int_0^t e^{-u}\mathbb E C_0(u)\,du\right)
 \le K e^t
\]
for a finite $K$.  It follows that
\[
 \mathbb E\langle\mathcal M\rangle_\infty
 =\int_0^\infty e^{-2u}
       \bigl(\mathbb E L(u)+\omega\mathbb E C_0(u)\bigr)\,du<\infty.
\]
Thus both terms on the right of \eqref{van:leaf-martingale} converge
almost surely and $e^{-t}L(t)\to V<\infty$.
From \eqref{van:leaf-rate}, $L$ admits a pathwise lower coupling
with a rate-one Yule process started with $L(0)$ particles.
Its $e^{-t}$-normalized limit is strictly positive almost surely.
Therefore $V>0$.

Since $\alpha<1$, \eqref{van:seed-amplitude} implies
$e^{-t}C(t)\to0$, whence $e^{-t}N(t)\to V$.
For the hitting time $\tau_n$ of population $n$, nonexplosion and
the Yule lower bound imply $\tau_n\to\infty$.
Consequently
\[
 \frac{C_n}{n^\alpha}
 =\frac{e^{-\alpha\tau_n}C(\tau_n)}
        {(e^{-\tau_n}N(\tau_n))^\alpha}
 \longrightarrow A V^{-\alpha}\in(0,\infty)
 \quad\text{almost surely}.
\]
A two-vertex seed is covered by applying the argument after its
first event and absorbing the finite time shift into $A$ and $V$.

\subsection{Effective computation of the root}

The transform admits the explicit convergent series
\begin{equation}
 F(s)=e^{-1}\sum_{j=0}^\infty\frac1{j!(s+j+1)},
 \qquad s>-1.
 \label{van:series}
\end{equation}
For $0\le s\le1$, truncation after $j=J$ has error at most
\[
 \frac{2}{(J+2)(J+1)!}.
\]
This follows by bounding $s+j+1\ge J+2$ in the tail and using
$\sum_{j\ge J+1}1/j!\le2/(J+1)!$.
Rational enclosures of $e^{-1}$ and a computable $\omega$ therefore
give certified enclosures of $m_\omega(s)$ at positive rational $s$.

The endpoint limits proved above allow effective discovery of rational
$a,b$ with $0<a<\alpha_\omega<b<1$ and
$m_\omega(a)>1>m_\omega(b)$.
Strict monotonicity gives certified root isolation.  To avoid relying
on deciding equality at a trial point, use the slope bound
\eqref{van:slope}: if a computed enclosure proves
$|m_\omega(c)-1|\le\delta$ for $c\in(0,1)$, then
$|\alpha_\omega-c|\le4\delta$.
Interval evaluation and this estimate supply terminating rational
enclosures of any requested width.

\begin{proof}[Completion of the proof of Theorem~\ref{van:main}]
The excursion calculation and renewal identity establish
\eqref{van:m} and \eqref{van:root}.  The first-child and endpoint
estimates prove uniqueness, the location $(0,1)$ and
\eqref{van:slope}.  The CMJ theorem, positive-event inheritance,
finite-family decomposition and leaf-clock argument give all
almost-sure limits in \eqref{van:limits}.  The last subsection
proves effectivity.
\end{proof}

\begin{remark}[Model distinction]
\label{van:scope}
For the original isotropic-redirection friend tree, a core may
redirect to another core.  Neighbouring environments then influence
its future leaf count, and the independent reproduction construction
\eqref{van:life} no longer holds.  The theorem supplies an exact
characterization for the modified VAN family only.  The original friend-tree exponent is not identified by this
characteristic equation. We use no discrete-time $L^1$ or
expected-amplitude conclusion.
\end{remark}

